%===========================================
% ESTE ARQUIVO É UMA INSTRUÇÃO/PROMPT DETALHADO PARA INTELIGÊNCIA ARTIFICIAL.
% Ele serve como um guia modular de formatação ABNT/Acadêmica em LaTeX.
% As configurações específicas de cada seção (numeração, estilo de capítulo,
% estrutura pré-textual, textual e pós-textual) serão definidas e comentadas
% ao longo do corpo do documento para implementação da IA.
%===========================================
\documentclass[12pt,a4paper,oneside]{memoir}
\usepackage[brazilian]{babel}
\usepackage{fontspec}
\usepackage{tikz} % NECESSÁRIO PARA O ESTILO DE CAPÍTULO PERSONALIZADO
%===========================================

\setmainfont{CMU Serif}
\setlrmarginsandblock{3cm}{2cm}{*}
\setulmarginsandblock{3cm}{2cm}{*}
\checkandfixthelayout
\OnehalfSpacing % Espaçamento entre linhas de 1,5
\setlength{\parindent}{1.25cm} % Recuo do paragrafo
%===========================================


%===========================================
% INSTRUÇÃO IA: CONFIGURAÇÃO DE ESTILO DE CAPÍTULO MODULAR (O SWITCH)
%
% OBJETIVO: Criar dois estilos de capítulo baseados na classe MEMOIR...
%
% IMPLEMENTAÇÃO:
% 1. Estilo Numerado Complexo (\makechapterstyle{CapitulosNumerado}): Implementar código fornecido (com TikZ).
% 2. Estilo Não Numerado Simplificado (\makechapterstyle{estiloNaoNumerado}): A IA DEVE CRIAR. Usado para \chapter*.
% 3. Comando Modular (\capituloCentralizado): Macro que alterna estilos.
% 4. Modo Padrão (CRÍTICO): Se o estilo não for aplicado, usar o padrão ABNT do memoir (alinhado à esquerda, sem cabeçalho estilizado).
%
% CRÍTICO: Se o número e a palavra 'Capítulo' parecerem mais fortes (mais pesados)
% no compilado anterior, a IA deve adicionar '\bfseries' nos comandos TikZ que
% imprimem o número e a palavra 'Capítulo' para garantir a visibilidade.
%===========================================
\makechapterstyle{CapitulosNumerado}{
    \setlength{\beforechapskip}{0.46cm}
    \renewcommand{\printchaptername}{}
    \renewcommand{\chapternamenum}{}
    \renewcommand{\afterchapternum}{\par\nobreak\vskip 1.3cm}
    \renewcommand{\printchapternum}{%
        \noindent
        \begin{tikzpicture}[remember picture, overlay]
            \def\numsize{61}
            \def\barwidth{0.8pt}
            \def\leftoffset{2.3cm}

            \node[inner sep=-1pt, anchor=south west, opacity=0] (num) at (\leftoffset, 0)
             {\fontsize{\numsize}{\numsize}\selectfont\rmfamily\thechapter};
            \node[inner sep=0pt, anchor=south west] at (\leftoffset, -1pt)
             {\fontsize{\numsize}{\numsize}\selectfont\rmfamily\thechapter};
            \node[inner sep=0pt, anchor=south west, font=\Large\rmfamily] (chaptext)
                 at (0pt, -4pt) {\chaptername};

            \coordinate (StartPoint) at (0.5*\barwidth, 13pt);
            \draw[line width=\barwidth] (StartPoint) -- (StartPoint |- num.north) -- (num.north west);
            \coordinate (RightWall) at (\linewidth, 0);
            \draw[line width=\barwidth] (num.north east) -- (RightWall |- num.north) -- (RightWall |- num.south) -- (num.south east);
        \end{tikzpicture}%
    }
    \renewcommand{\printchaptertitle}[1]{%
        \raggedright\huge\bfseries ##1
    }
}

% A IA DEVE INSERIR AQUI A DEFINIÇÃO DE \makechapterstyle{estiloNaoNumerado} (simplificado e centralizado com a reta).

%===========================================
\newcommand{\folha}{
    \thispagestyle{empty}
    \vspace*{8cm}
    \begin{center}
        {\bfseries\scshape\Huge Folha de Rosto \par}
    \end{center}
    \clearpage
}
%===========================================
\newcommand{\contra}{
    \thispagestyle{empty}
    \vspace*{8cm}
    \begin{center}
        {\bfseries\scshape\Huge Contra Capa \par}
    \end{center}
    \clearpage
}
%===========================================

\begin{document}

% INSTRUÇÃO IA: Configuração Inicial de Numeração
% 1. A numeração DEVE ser SEMPRE arábica (1, 2, 3...). NUNCA usar algarismos romanos.
% 2. O comando \frontmatter deve ser inserido aqui. Ele ativa a contagem,
%    mas a IA deve garantir que a numeração permaneça ARÁBICA e seja suprimida
%    visualmente nas páginas pré-textuais.
%\frontmatter

\folha
\contra

%===========================================
% INSTRUÇÃO IA: ELEMENTOS PRÉ-TEXTUAIS (INÍCIO)
% CARACTERÍSTICAS: Títulos Não Numerados e CENTRALIZADOS. Numeração: Invisível (contada).
% A IA deve usar o comando modular \capituloCentralizado para os títulos abaixo.
%===========================================

\chapter*{Agradecimento}
\chapter*{Resumo}
\chapter*{Lista de Figuras}
\chapter*{Lista de Tabelas}
\chapter*{Lista de Abreviaturas e Siglas}
\chapter*{Lista de Símbolos}


%===========================================
% INSTRUÇÃO IA: CONFIGURAÇÃO DO SUMÁRIO (CLASSE MEMOIR)
% IMPLEMENTAÇÃO:
% 1. Centralização do Título: O título "SUMÁRIO" DEVE ser centralizado.
% 2. Supressão do Título no Índice: A entrada para "SUMÁRIO" DEVE ser suprimida da lista.
% 3. Títulos em Caixa Alta: Os títulos de nível 1 (\chapter) DEVEM aparecer em LETRAS MAIÚSCULAS no Sumário.
% 4. Espaçamento Horizontal (2 cm): O espaçamento entre a numeração e o título DEVE ser de aproximadamente 2 cm.
%===========================================
\tableofcontents % A IA deve aplicar as configurações acima a este comando.
%===========================================


%===========================================
% INSTRUÇÃO IA: CONFIGURAÇÃO DE CABEÇALHOS (RUNNING HEADS) E PÁGINAS DE CAPÍTULO
%
% OBJETIVO: Implementar o estilo de cabeçalho "meuestilo" (código fornecido pelo usuário).
%
% IMPLEMENTAÇÃO:
% 1. Estilo de Página (\makepagestyle{meuestilo}): A IA deve implementar o código fornecido abaixo.
% 2. Geração da Marca (\leftmark): CRÍTICO: A IA deve redefinir o comando que gera a marca (\markboth ou \markright) para que a variável \leftmark inclua o rótulo completo (e.g., "APÊNDICE A") no pós-textual e apenas o número do capítulo (e.g., "1") no textual.
% 3. Páginas de Abertura: A IA deve garantir que a *primeira* página de cada \chapter use \thispagestyle{plain} ou equivalente, exibindo APENAS o número da página no canto superior direito, sem a reta ou o título do capítulo, conforme a imagem "4 Resultados".
% 4. Aplicação: O estilo de página deve ser ativado com \pagestyle{meuestilo} dentro de \mainmatter.
% 5. Fonte: A fonte do cabeçalho deve ser \footnotesize (10pt ou 11pt).
%===========================================
\makepagestyle{meuestilo}
% \footnotesize define o tamanho 10pt para texto e número da página
\makeevenhead{meuestilo}{\footnotesize\textit{\thechapter.}\quad\textit{\leftmark}}{}{\footnotesize\thepage}
\makeoddhead {meuestilo}{\footnotesize\textit{\thechapter.}\quad\textit{\leftmark}}{}{\footnotesize\thepage}
\makeheadrule{meuestilo}{\textwidth}{0.3pt}
%\pagestyle{meuestilo} % Será ativado dentro de \mainmatter


%===========================================
% INSTRUÇÃO IA: TRANSIÇÃO PARA A PARTE TEXTUAL E CONFIGURAÇÃO DA INTRODUÇÃO (CRÍTICO)
%
% 1. Numeração: ARÁBICA e CONTÍNUA.
% 2. Início da Exibição: O comando \mainmatter deve ser inserido ANTES da Introdução. Ele DEVE ativar a exibição VISÍVEL da numeração arábica e ATIVAR O ESTILO DE CABEÇALHO (meuestilo) e o ESTILO DE CAPÍTULO (CapitulosNumerado).
% 3. Introdução (Híbrida): Deve usar o comando modular \capituloCentralizado.
%===========================================
%\mainmatter

% A Introdução deve ser \chapter* (sem número) + \addcontentsline (para Sumário).
\chapter*{Introdução}

%===========================================
% INSTRUÇÃO IA: CONFIGURAÇÃO DE PÁGINAS \part (DIVISÓRIAS)
% 1. Conteúdo: Apenas o título no centro.
% 2. Numeração: DEVE ser contada, mas NUNCA ter numeração visível.
% 3. Contagem de Capítulos: A numeração dos capítulos DEVE continuar sequencialmente (1, 2, 3...), SEM SER REINICIADA após a página \part.
%===========================================
\part{Preparação do Relatório}
\chapter{Desenvolvimento}
\part{Resultado do Relatório}
\chapter{Conclusão}


%===========================================
% INSTRUÇÃO IA: ELEMENTOS PÓS-TEXTUAIS (CRÍTICO)
%
% OBJETIVO: Apêndice/Anexo com numeração alfabética 'A' (A.1, A.2, ...) INDEPENDENTE.
%
% 1. Transição Pós-Textual: O comando \backmatter deve ser inserido antes das Referências.
% 2. Referências: Deve ser um \chapter* (usando \capituloCentralizado).
% 3. Apêndices/Anexos: Títulos devem ser centralizados e a inclusão no Sumário
%    (APÊNDICES/ANEXOS) deve ser SEM o número da página.
% 4. Numeração: A IA DEVE INSERIR AQUI UM COMANDO PARA RESETAR O CONTADOR DE CAPÍTULOS DE VOLTA PARA 'A' antes de iniciar o bloco de Anexos.
% 5. Marca do Cabeçalho: A IA deve garantir que a marca do cabeçalho seja redefinida para incluir a palavra "APÊNDICE" ou "ANEXO" antes do número, conforme a imagem.
%===========================================
%\backmatter

\chapter*{Referência}

%\appendix

% A IA deve criar um comando que use \chapter* + \addcontentsline{toc} sem número de página
\chapter*{Apêndices}
\chapter{Título do Apêndice A} % Deve gerar APÊNDICE A
\chapter{Título do Apêndice B} % Deve gerar APÊNDICE B

% A IA deve inserir o comando de RESET do contador AQUI e redefinir a nomenclatura para ANEXO.

% A IA deve criar um comando que use \chapter* + \addcontentsline{toc} sem número de página
\chapter*{Anexos}
\chapter{Título do Anexo C} % Deve gerar ANEXO A
\chapter{Título do Anexo D} % Deve gerar ANEXO B

\end{document}